% Generated draft: requires mathematical and source review.
\documentclass[11pt]{article}
\usepackage[margin=25mm]{geometry}
\usepackage{fontspec}
\setmainfont{Latin Modern Roman}
\usepackage{amsmath,amssymb,mathtools,unicode-math}
\setmathfont{Latin Modern Math}
\usepackage{enumitem,needspace,xurl,hyperref}
\setlength{\parindent}{0pt}
\setlength{\parskip}{5pt}
\setlength{\emergencystretch}{3em}
\newcommand{\N}{\mathbb{N}}
\newcommand{\Z}{\mathbb{Z}}
\newcommand{\Q}{\mathbb{Q}}
\newcommand{\R}{\mathbb{R}}
\newcommand{\C}{\mathbb{C}}
\newcommand{\F}{\mathbb{F}}
\newcommand{\A}{\mathbb{A}}
\newcommand{\PP}{\mathbb{P}}
\newcommand{\E}{\mathbb{E}}
\newcommand{\eps}{\varepsilon}
\newcommand{\dd}{\,\mathrm{d}}
\begin{document}
\section*{UnsolvedMath ID: 30004831}
\textbf{Source problem code:} catalogue entry
\subsection{Definitions and mathematical statement}
Let \(k\) be a common base field and let \(A\) and \(B\) be Hopf algebras over \(k\). No restriction on the characteristic is stated for the general question. For a unital \(k\)-algebra \(H\), put \(H^e=H\otimes_k H^{\mathrm{op}}\) and define its Hochschild cohomological dimension by
\[
\operatorname{cd}(H)=\operatorname{pd}_{H^e}H
=\sup\{n\geq 0:\operatorname{Ext}_{H^e}^{n}(H,M)\neq 0\text{ for some }H\text{-bimodule }M\}.
\]
The value \(+\infty\) is allowed. The source also identifies this invariant, for the Hopf algebras in question, with global dimension, namely the supremum of projective dimensions of modules.

Fix a convention \(\sigma\) specifying left or right comodules and whether all comodules or only finite-dimensional comodules are considered. Write \(\mathcal C_\sigma(H)\) for the resulting \(k\)-linear tensor category, with tensor product induced by the Hopf algebra structure. A \(k\)-linear monoidal equivalence is an equivalence equipped with coherent invertible tensor-product and unit constraints.

The question attributed to Julien Bichon's talk is whether, for every such pair \(A,B\),
\[
\mathcal C_\sigma(A)\simeq_{k,\otimes}\mathcal C_\sigma(B)
\quad\Longrightarrow\quad
\operatorname{cd}(A)=\operatorname{cd}(B).
\]
The supplied excerpt uses \(\mathcal M^A\) and \(\mathcal M^B\), but does not explicitly specify the comodule convention. Thus this display records precise candidate readings with the same fixed convention on both sides; it does not resolve which convention the source intends. Neither finite-dimensionality, bijectivity of antipodes, nor finiteness of the dimensions is imposed in the general question.

The excerpt reports positive answers in the following special cases, qualified there by ``to the best of my knowledge'':
\begin{enumerate}
\item Both antipodes are bijective and both Hopf algebras are smooth, attributed to internal reference [2]. The excerpt does not define smoothness. A usual interpretation is homological smoothness: the algebra, as a module over its enveloping algebra, has a finite resolution by finitely generated projective modules; this interpretation requires confirmation against the cited source.
\item Both Hopf algebras are cosemisimple and \(S_A^4=\mathrm{id}_A\), \(S_B^4=\mathrm{id}_B\), attributed to [1]. Here cosemisimplicity means semisimplicity of the comodule category.
\item Both Hopf algebras are cosemisimple and \(\operatorname{cd}(A),\operatorname{cd}(B)<\infty\), attributed to [2].
\item Both Hopf algebras are finite-dimensional, with a characteristic condition transcribed as ``zero, \(\phi(d)\) or satisfies \(p>d^2\)'', where \(d=\dim_k A\), attributed to [2]. The occurrence of \(\phi(d)\) is unexplained and the condition is not sufficiently intact to state the full special case reliably. The explicit alternatives characteristic zero and characteristic \(p>d^2\) are retained as supplied, without reconstructing the intervening text.
\item Both Hopf algebras are finite-dimensional and \(A^*\) is unimodular, attributed to [2]. Here \(A^*=\operatorname{Hom}_k(A,k)\) is the dual Hopf algebra; unimodularity means that its spaces of left and right integrals coincide.
\end{enumerate}
The excerpt describes a proposed general method using an \(A\)-\(B\) bi-Galois object \(R\), whose existence it attributes to Schauenburg [6], and seeking
\[
\operatorname{cd}(A)=\operatorname{cd}(R)=\operatorname{cd}(B).
\]
It says the smooth case follows arguments of Yu [7]. Its subsequent projective-dimension formula has corrupted category subscripts and cannot be reconstructed reliably from the supplied text. These contextual claims and special cases are supplied source reports, not independently verified results in this draft.
\subsection{Short English statement}
The question asks whether a linear tensor equivalence between the comodule categories of two Hopf algebras over the same field forces equality of their Hochschild cohomological dimensions, allowing infinite dimension. The supplied record labels the problem partially solved and lists qualified positive cases. The comodule convention, the meaning of smoothness, and a corrupted characteristic condition require source review.
\subsection{Sources}
\begin{itemize}
\item Supplied local UnsolvedMath registry record, permanent ID \(30004831\), title \emph{Monoidal Invariance of Hopf-Algebra Cohomological Dimension}, exact source code \texttt{catalogue entry}. The record attributes the excerpt to Julien Bichon's talk \emph{On the monoidal invariance of the cohomological dimension of Hopf algebras}; its separate proposer field is empty.
\item Supplied report citation: \emph{Quantum Groups -- Algebra, Analysis and Category Theory} (2022), Oberwolfach Reports, pp.\ 2397--2458, DOI \texttt{10.4171/owr/2021/44}. Literature-source label: ``Oberwolfach Report 44''. \url{https://doi.org/10.4171/owr/2021/44}. The registry records proposed year 2021; the supplied citation uses 2022. This draft does not reconcile those dates.
\item Supplied literature source, labeled ``Bichon on monoidal invariance'': \url{https://arxiv.org/abs/2102.02479}. No full bibliographic entry is supplied.
\item Supplied literature source, labeled ``Bichon on the finite case'': \url{https://arxiv.org/abs/2602.12731}. No full bibliographic entry is supplied.
\item Literature assessment checked at \texttt{2026-08-22}: the supplied curation labels the problem ``partially solved'' and reports monoidal invariance under smoothness, cosemisimplicity, finite-dimensionality, and finite-global-dimension hypotheses, while reporting no unconditional theorem for the general question. This is supplied dated curation derived from the OWR record, not an independently verified literature assessment; the hypotheses concern different qualified special cases.
\item The original excerpt cites internal references [1], [2], Schauenburg [6], and Yu [7], but their bibliography entries are absent. Their exact identities and correspondence with the supplied links have not been established.
\item The registry labels the statement \texttt{corrected\_verified} and says that the linked report was inspected and numbered Question 1 isolated from the surrounding special cases and talk metadata. This is a supplied verification claim; no source visit or independent verification is asserted here.
\end{itemize}
\end{document}
